\documentclass{article}

\usepackage{amssymb}
\usepackage{mathtools}
\usepackage{amsmath}
\usepackage{amsthm}
\usepackage{babel}
\usepackage{stackengine}
\usepackage{tikz}
% Smooth peak profile for plotting radial deformations.
% The small cutoff avoids division by zero at peak boundaries.
\pgfmathdeclarefunction{smoothpeak}{1}{%
	\pgfmathparse{
		exp(1-1/max(0.01,sin(#1)*sin(#1)))
	}%
}
\usepackage{tikz-cd}
\usepackage{quiver}
\usepackage{mathrsfs}
\usepackage{braket}
\usepackage{float}
\usepackage{scalerel}
\usepackage{stackengine,wasysym}
\usepackage{blindtext}
\usepackage[final]{hyperref}
\usepackage{bbold}
\usepackage{geometry}
\usepackage{cancel}
\usepackage{booktabs}
\usepackage{comment}

\title{Geometric Asymmetry and Gravitational Fields of Closed Elastic Loops}
\author{Ryan J. Buchanan\footnote{Active Inference Institute. \href{https://orcid.org/0009-0000-1253-0031}{ORCiD: 0009-0000-1253-0031}}}
\date{October, 2026}

\newtheorem{war}{Warning}
\newtheorem{dn}{Definition}
\newtheorem{as}{Assumption}
\newtheorem{rk}{Remark}
\newtheorem{eg}{Example}
\newtheorem{pp}{Proposition}
\newtheorem{ax}{Axiom}
\newtheorem{tm}{Theorem}
\newtheorem{lem}{Lemma}
\newtheorem{nt}{Notation}
\newtheorem{prop}{Property}
\newtheorem{rc}{Recollection}
\newtheorem{exc}{Exercise}


\numberwithin{tm}{section}
\numberwithin{dn}{section}
\numberwithin{as}{section}
\numberwithin{rk}{section}
\numberwithin{lem}{section}
\numberwithin{ax}{section}
\numberwithin{pp}{section}
\numberwithin{eg}{section}
\numberwithin{war}{section}
\numberwithin{equation}{section}
\numberwithin{prop}{section}
\numberwithin{nt}{section}
\numberwithin{rc}{section}


\newcommand\restr[2]{{% we make the whole thing an ordinary symbol
		\left.\kern-\nulldelimiterspace % automatically resize the bar with \right
		#1 % the function
		\littletaller % pretend it's a little taller at normal size
		\right|_{#2} % this is the delimiter
}}

\def\code#1{\texttt{#1}}


\begin{document}
\maketitle


\begin{abstract}
	We develop a geometric and variational framework for investigating
	asymmetric closed elastic loops and their associated gravitational
	fields. Modeling a loop as a smooth closed curve in Euclidean
	three-space, we introduce an elastic energy combining line tension
	and bending rigidity, and consider families of deformations that
	allow the geometric asymmetry of the loop to be varied. An
	effective line density derived from the elastic energy provides a
	phenomenological source for a weak-field Newtonian potential, whose
	far-field structure can be analyzed through multipole expansions.
	We emphasize the distinction between this effective description
	and a gravitational field derived from the full stress-energy tensor
	of a relativistic string.
	
	We further model branes as embedded manifolds equipped with
	admissible ambient-diffeomorphism flows. A joint energy functional
	combining string, brane, and interaction energies defines a
	variational framework for selecting compatible brane configurations
	and motivates a coupled gradient-flow model for their dissipative
	evolution. Finally, we discuss a controlled extension to immersed
	loops with transverse self-intersections, identifying additional
	issues involving gravitational self-interaction and topology-changing
	reconnection. The resulting framework provides a mathematical
	setting for studying geometric asymmetry, effective gravitational
	fields, and coupled string--brane deformation, without presupposing
	a fundamental physical interpretation of the model.
\end{abstract}

\pagebreak

\tableofcontents

\section{Introduction}
Loop spaces arestandard objects in algebraic topology.
While the theory of infinite loop spaces and its connections with
stable homotopy theory are initially developed, for example, in
\cite{Adams}, loop-space methods also arise in mathematical physics.
Schreiber studies their role in supersymmetric quantum
mechanics, superstring background fields, and nonabelian
strings \cite{Schreiber}.

Let $(X,x_0)$ be a pointed topological space, and let
$s_0\in S^1$ be a chosen basepoint. The based loop space
of $X$ is
\begin{equation}
	\Omega_{x_0}X
	\coloneq
	\operatorname{Map}_*
	\bigl((S^1,s_0),(X,x_0)\bigr).
\end{equation}
The free loop space is
\begin{equation}
	LX
	\coloneq
	\operatorname{Map}(S^1,X).
\end{equation}

At a fixed worldsheet time, a closed-string configuration
can be represented kinematically by a map
$\gamma:S^1\to X$, whereas an open-string configuration
can be represented by a map $\gamma:[0,1]\to X$,
together with appropriate endpoint conditions. Thus, the
free loop space provides a natural configuration space for
closed strings. A based loop space imposes the additional
condition that a chosen point of the circle maps to $x_0$.\footnote{A natural question we ask is when exactly it is appropriate to use the latter construction; for instance, if and when we have a metric which varies along the string, or in other cases where $x_0$ has a manifest physical interpretation.} 

In perturbative string theory, the closed-string spectrum
contains a massless spin-two state interpreted as the
graviton \cite{Green,Polchinski,Zwiebach}. This fact
motivates our investigation of geometric loop
configurations and their possible gravitational effects.
However, a loop configuration and a graviton state are
distinct objects, and their relationship requires a
separate physical model.

In this paper, we begin with a geometric model of closed
elastic loops.\footnote{Investiagions into elastic curves, or ``cables" as they were called in \cite{Atai} are well-documented.} We investigate how their shape, curvature,
and energy distribution vary under deformation, and how
these quantities may be used to calculate their
gravitational fields in a specified approximation.

Strings are not the only extended objects considered in
string theory. D-branes provide boundary conditions for
open strings, while more general $p$-branes are extended
objects with $p$ spatial dimensions. For the present geometric model, let $X$ be a smooth
Riemannian manifold and let an $n$-brane be represented
at each stage of its evolution by an embedded
$n$-dimensional submanifold $Y_t\subset X$. We regard
the brane as a dynamical geometric object and describe
its admissible deformations through a specified
collection of ambient diffeomorphism flows.


The choice of ambient diffeomorphisms is motivated by two considerations. First, a smooth diffeomorphism of the ambient manifold transports the brane without changing its diffeomorphism type. In particular, a compact genus-zero brane remains genus zero throughout any admissible flow. Second, the smoothness of the deformation is essential for the geometric analysis that follows: the quantities of interest, including curvature and bending energy, depend on differential information that an arbitrary homeomorphism need not preserve. The collection of admissible flows therefore serves to specify which geometric deformations are permitted within the model. Requiring the flows to be defined on a common time interval, with the identity transformation as their initial condition, allows each evolution to be understood relative to a fixed reference configuration. Different choices of admissible flows may impose different restrictions on the deformation of the brane, and hence determine which configurations can be reached from the initial embedding. 
\begin{war}
	At this stage, we do not prescribe a particular equation of motion or require the brane to follow an energy-minimizing trajectory. The flow set instead records the admissible geometric evolutions independently of the dynamical law that may subsequently select among them. This distinction allows us to introduce a metric-dependent energy functional later and investigate its critical points and induced evolution without building the desired dynamics into the definition of a brane from the outset.
\end{war}

Since we are building flow sets into the definition of a brane (Def.~\ref{brane}), the geometry of allowable deformations is strictly controlled from the outset. So, given a $\Phi_t \in \mathscr{F}_\mathfrak{B}$, there is guarunteed to be at least one brane geometry $Y_t$ for our dynamics. However, we also wish to address the additional deformations of the string $\gamma$ itself, which possesses its own geometric degrees of freedom. We represent a closed-string configuration by a regular loop $\gamma:S^1\to X$ and allow this loop to deform independently of the brane. In particular, the brane flow does not determine the string's shape, curvature, or distribution of elastic energy. These quantities must be accounted for by a separate family of loop configurations. Consequently, the configurations of interest consist of a string and a brane whose respective deformations
are governed by distinct admissibility conditions. Their interaction is then determined by a common
geometric functional, which couples the string's elastic energy to the evolving brane geometry.

Let us restrict attention to closed strings without
self-intersections. Accordingly, each string configuration is
represented by a smooth embedding
\[
\gamma:S^1\hookrightarrow X.
\]
This condition ensures that the string is a regular embedded
curve, so its geometric quantities, including arclength and
curvature, are well-defined along the entire loop. We therefore
take the initial configuration space of strings to be
$\operatorname{Emb}(S^1,X)$, rather than the larger space of
immersed loops.

This restriction is imposed purely for mathematical convenience; self-intersecting strings are reserved for a later extension of the theory  in \S\ref{sec:self-intersections}. There, we may enlarge the configuration space to include immersed loops with prescribed classes of singularities, such as finitely many transverse double points,
subject to suitable regularity conditions. Such an extension will require us to examine how self-intersections affect the energy functional, the admissible deformations, and the resulting gravitational field.

\section{Geometric Setup and Elastic-Loop Model}
\label{sec:geometric-model}

\subsection{Configuration Space of Strings and Branes}
\label{subsec:configuration-space}

Let $(X,g)$ be a smooth Riemannian manifold. We initially
restrict attention to closed strings represented by smooth
embeddings
\[
\gamma:S^1\hookrightarrow X,
\]
thereby excluding self-intersections from the initial model.

We regard a brane as an embedded submanifold whose admissible
geometric deformations are specified by a collection of
ambient diffeomorphism flows. This approach is inspired by
the theory of diffeomorphic surface flows developed in
\cite{ZhangEtAl}.

\begin{dn}[Brane] \label{brane}
	Let $X$ be a smooth manifold and let $\Sigma^n$ be a smooth
	$n$-dimensional manifold. An $n$-brane $\mathfrak{B}^n$ in $X$
	is a triple
	\[
	\mathfrak{B}^n
	=
	(\Sigma^n,\iota_0,\mathscr{F}_{\mathfrak{B}}),
	\]
	where $\iota_0:\Sigma^n\hookrightarrow X$ is a smooth
	embedding and $\mathscr{F}_{\mathfrak{B}}$ is a specified
	collection of smooth flows of ambient diffeomorphisms
	$\Phi_t:X\to X$, defined on a common time interval $I$
	containing $0$, with $\Phi_0=\operatorname{id}_X$.
	
	Each flow induces a family of embeddings
	\[
	\iota_t=\Phi_t\circ\iota_0,
	\qquad
	Y_t=\iota_t(\Sigma^n)\subset X.
	\]
	The collection $\mathscr{F}_{\mathfrak{B}}$ is called the
	\emph{flow set} of $\mathfrak{B}^n$ and specifies its
	admissible geometric deformations. Every such deformation
	preserves the diffeomorphism type of $\Sigma^n$.
\end{dn}

We now consider the relationship between the brane's
admissible deformations and the geometry of a closed string.
Let
\[
q=(\gamma,\iota)
\]
denote a configuration consisting of a smooth embedded loop
$\gamma:S^1\hookrightarrow X$ and a brane embedding
$\iota:\Sigma^n\hookrightarrow X$.

The flow set $\mathscr{F}_{\mathfrak{B}}$ determines the
admissible geometric realizations of the brane. Writing
$Y_0=\iota_0(\Sigma^n)$, an ambient diffeomorphism
$\Phi_t:X\to X$ induces
\[
\iota_t=\Phi_t\circ\iota_0,
\qquad
Y_t=\Phi_t(Y_0).
\]
The resulting family remains within the ambient-diffeomorphism
orbit of the initial brane, and hence preserves its
diffeomorphism type.

The admissibility of a brane configuration does not, however,
determine whether it is geometrically compatible with a given
string. To express this dependence, we consider a joint energy
functional
\[
\mathcal{J}[\gamma,\iota]
=
\mathcal{E}_{\mathrm{str}}[\gamma]
+\mathcal{E}_{\mathrm{br}}[\iota]
+\mathcal{E}_{\mathrm{int}}[\gamma,\iota].
\]
The interaction term couples the string's geometry to the
current brane configuration. For a fixed string, this
functional provides a criterion for comparing admissible
brane embeddings. In particular, we may seek configurations
that minimize the brane and interaction energies over the
orbit generated by $\mathscr{F}_{\mathfrak{B}}$, whenever
such minimizers exist.

The string and brane are permitted to deform independently.
This independence is necessary if their relative position
and incidence are to vary during an evolution. Their motions
may nevertheless remain coupled through $\mathcal{J}$.
We distinguish, accordingly, between the flow set, which
specifies admissible brane deformations, and the dynamical
law, which determines how the string and brane evolve
together. Our master equation for the deformation of $\mathfrak{B}$ is Eq.~\eqref{Master}.

\subsection{A Four-Parameter Family of Closed Loops}
\label{subsec:loop-family}

For the initial geometric model, we take $X=\mathbb{R}^3$
and restrict the string to a fixed plane. We normalize the reference loop to be the unit circle.
Divide its angular parameter domain into two hemispheres,
$I_1=[0,\pi]$ and $I_2=[\pi,2\pi]$. We introduce four
parameters
\[
(A_1,N_1,A_2,N_2),
\]
where $A_i\geq 0$ controls the peak height on hemisphere $i$,
and $N_i\in\mathbb{N}$ controls its number of peaks.

To obtain a smooth family of curves, define the periodic
profile
\begin{equation}
	\psi(u)=
	\begin{cases}
		\exp\!\left(1-\dfrac{1}{\sin^2 u}\right),
		& \sin u\neq 0,\\[6pt]
		0,  & \sin u=0.
	\end{cases}
\end{equation}
This profile is smooth, has maximum value $1$, and vanishes
to all orders at integer multiples of $\pi$.

Define the radial function
\[
r(\theta)=
\begin{cases}
	1+A_1\psi(N_1\theta),
	& 0\leq\theta\leq\pi,\\[4pt]
	1+A_2\psi(N_2(\theta-\pi)),
	& \pi\leq\theta\leq 2\pi.
\end{cases}
\]
The corresponding closed curve is
\[
\gamma(\theta)=
\bigl(r(\theta)\cos\theta,\,
r(\theta)\sin\theta,\,0\bigr),
\qquad 0\leq\theta\leq 2\pi.
\]
The maximum radius on hemisphere $i$ is $1+A_i$, and the
angular spacing between adjacent peaks is
\[
\Delta_i=\frac{\pi}{N_i}.
\]

\begin{figure}[htbp]
	\centering
	\begin{tikzpicture}[
		scale=1.2,
		line cap=round,
		line join=round
		]
		
		% Reference unit circle
		\draw[
		gray!65,
		densely dashed,
		line width=0.7pt
		]
		(0,0) circle (1);
		
		% Coordinate axes
		\draw[gray!25, thin]
		(-2.05,0) -- (2.05,0);
		\draw[gray!25, thin]
		(0,-1.65) -- (0,1.65);
		
		% Upper hemisphere: A_1 = 0.8, N_1 = 2
		% Radial profile: r(theta) = 1 + A_1 psi(N_1 theta)
		\draw[
		orange!85!black,
		line width=1.3pt
		]
		plot[
		variable=\t,
		domain=0:180,
		samples=721
		]
		(
		{(1+0.8*smoothpeak(2*\t))*cos(\t)},
		{(1+0.8*smoothpeak(2*\t))*sin(\t)}
		);
		
		% Lower hemisphere: A_2 = 0.45, N_2 = 4
		% Radial profile: r(theta) = 1 + A_2 psi(N_2(theta-pi))
		\draw[
		blue!70!black,
		line width=1.3pt
		]
		plot[
		variable=\t,
		domain=180:360,
		samples=961
		]
		(
		{(1+0.45*smoothpeak(4*(\t-180)))*cos(\t)},
		{(1+0.45*smoothpeak(4*(\t-180)))*sin(\t)}
		);
		
		% Mark the upper peak locations (radius 1.8)
		\foreach \ang in {45,135} {
			\fill[orange!85!black]
			({1.8*cos(\ang)},{1.8*sin(\ang)})
			circle (1.7pt);
		}
		
		% Mark the lower peak locations (radius 1.45)
		\foreach \ang in {202.5,247.5,292.5,337.5} {
			\fill[blue!70!black]
			({1.45*cos(\ang)},{1.45*sin(\ang)})
			circle (1.7pt);
		}
		
		% Reference radius
		\draw[gray!75, thin]
		(0,0) -- (1,0);
		
		\node[
		font=\scriptsize,
		fill=white,
		inner sep=1.5pt
		] at (0.48,0.13) {$r=1$};
		
		% Hemisphere labels
		\node[
		font=\small\bfseries,
		text=orange!85!black
		] at (0,2.05)
		{$A_1=0.8,\quad N_1=2$};
		
		\node[
		font=\small\bfseries,
		text=blue!70!black
		] at (0,-2.05)
		{$A_2=0.45,\quad N_2=4$};
		
		% Angular spacing annotations
		\node[
		font=\small,
		fill=white,
		inner sep=1.5pt
		] at (1.75,1.45)
		{$\Delta_1=\pi/2$};
		
		\node[
		font=\small,
		fill=white,
		inner sep=1.5pt
		] at (1.75,-1.65)
		{$\Delta_2=\pi/4$};
		
	\end{tikzpicture}
	
	\caption{An asymmetric closed loop with independently
		controlled peak heights and spacings on its two
		hemispheres. The dashed circle is the unit-circle
		reference configuration. Here $A_1=0.8$, $N_1=2$,
		$A_2=0.45$, and $N_2=4$, giving maximum radii
		$1+A_1=1.8$ and $1+A_2=1.45$, respectively.}
	
	\label{fig:asymmetric-closed-loop}
\end{figure}

Our motivation for splitting the ``graviton" (closed loop) into two hemispheres, is that gravitational ``charges" should act along an an axis, with the direction being deterimined by the asymmetry between the two halves. Our four-parameter setup is a simplified toy model, as we do not expect the deformations to be distributed so evenly.


\subsection{Geometric Energy and Energy Asymmetry}
\label{subsec:energy-asymmetry}

Let primes denote derivatives with respect to $\theta$.
The arclength element of the radial curve is
\[
ds=\sqrt{r^2+(r')^2}\,d\theta.
\]
Its curvature magnitude is
\[
\kappa(\theta)=
\frac{\left|r^2+2(r')^2-rr''\right|}
{\left(r^2+(r')^2\right)^{3/2}}.
\]

We consider the elastic-loop energy
\begin{equation}
	\label{eq:string-energy}
	E_{\mathrm{str}}[\gamma]
	=
	\int_\gamma
	\left(\tau+\frac{\kappa_b}{2}\kappa^2\right)ds,
\end{equation}
where $\tau>0$ is the line-tension coefficient and
$\kappa_b>0$ is the bending-rigidity coefficient.

For $i\in\{1,2\}$, let $E_i$ denote the energy contributed
by hemisphere $I_i$. Explicitly,
\begin{equation}
	\label{eq:hemisphere-energy}
	E_i
	=
	\int_{I_i}
	\left(\tau+\frac{\kappa_b}{2}\kappa(\theta)^2\right)
	\sqrt{r(\theta)^2+r'(\theta)^2}\,d\theta.
\end{equation}
The total energy is
\begin{equation}
	\label{eq:total-energy}
	E_{\mathrm{total}}=E_1+E_2.
\end{equation}

To quantify the unevenness of the energy distribution, we
define the normalized hemispheric energy asymmetry by
\begin{equation}
	\label{eq:energy-asymmetry}
	\mathcal{A}_E
	=
	\frac{E_2-E_1}{E_1+E_2},
\end{equation}
provided $E_1+E_2>0$. Thus $\mathcal{A}_E=0$ when the two
hemispheres contribute equal energy, while its sign records
which hemisphere contributes more. The dependence of
$E_{\mathrm{total}}$ and $\mathcal{A}_E$ on the four shape
parameters will be investigated numerically in \S\ref{subsec:numerical-example}.


\subsection{From Geometric Energy to Gravitational Fields}
\label{subsec:gravity-preview}

The preceding construction determines the loop geometry and
its modeled energy distribution. To investigate its
gravitational influence, we must additionally specify how
this energy sources a gravitational field.

In the weak-field approximation, we will first consider an
effective line-mass density derived from the chosen energy
density. We will then calculate the gravitational potential
and examine its multipole expansion, distinguishing the
contribution of total mass from the effects of the spatial
distribution of energy.

The gravitational model and its assumptions are developed
in Section~\ref{sec:multipoles}.

\section{Multipole Analysis}
\label{sec:multipoles}

We now investigate how the geometric energy distribution of a
closed elastic loop determines its gravitational field. Our
starting point is the effective energy functional
\eqref{eq:string-energy}. Since the loop is an extended source,
its gravitational influence depends not only on its total energy,
but also on the spatial distribution of that energy. Multipole
analysis provides a systematic means of separating these
contributions \cite{PoissonWill2014Gravity}.

Throughout this section, we work in a weak-field, quasi-static
approximation. We treat the energy density derived from
\eqref{eq:string-energy} as an effective source for a Newtonian
gravitational potential. This is a simplifying assumption of
our geometric model, rather than a complete relativistic
description of a string source.

\subsection{Effective Line-Mass Density}

Let $\gamma:S^1\to\mathbb{R}^3$ be a regular closed loop,
parametrized by arclength $s$, and let
$\kappa(s)$ denote its curvature. Define the geometric energy
density per unit length by
\begin{equation}
	\varepsilon_{\mathrm{str}}(s)
	\coloneq
	\tau+\frac{\kappa_b}{2}\kappa(s)^2.
\end{equation}
The total string energy is consequently
\begin{equation}
	E_{\mathrm{total}}
	=
	\oint_\gamma
	\varepsilon_{\mathrm{str}}(s)\,ds.
\end{equation}

For the purposes of the present approximation, we associate
to this energy density the effective line-mass density
\begin{equation}
	\label{eq:effective-line-mass}
	\lambda_{\mathrm{eff}}(s)
	\coloneq
	\frac{\varepsilon_{\mathrm{str}}(s)}{c^2},
\end{equation}
where $c$ denotes the speed of light. The corresponding
effective mass is
\begin{equation}
	\label{eq:effective-total-mass}
	M_{\mathrm{eff}}
	=
	\oint_\gamma \lambda_{\mathrm{eff}}(s)\,ds
	=
	\frac{E_{\mathrm{total}}}{c^2}.
\end{equation}

Equation~\eqref{eq:effective-line-mass} provides a convenient
bridge between the geometric energy model and Newtonian
gravity. A physically complete relativistic treatment would
require the stress-energy tensor of the string, including
its tension and any dynamical contributions
\cite{ChernoffFlanaganWardell2019StringBackreaction}.

\subsection{Newtonian Gravitational Potential}

Choose a fixed origin $O\in\mathbb{R}^3$ and write
$\mathbf{y}(s)$ for the position vector of the string relative
to this origin. In the effective Newtonian model, the
gravitational potential at an observation point
$\mathbf{x}\notin\gamma(S^1)$ is
\begin{equation}
	\label{eq:newtonian-loop-potential}
	\Phi(\mathbf{x})
	=
	-G\oint_\gamma
	\frac{\lambda_{\mathrm{eff}}(s)}
	{|\mathbf{x}-\mathbf{y}(s)|}\,ds,
\end{equation}
where $G$ is Newton's gravitational constant. The corresponding
gravitational acceleration is
\begin{equation}
	\mathbf{g}(\mathbf{x})=-\nabla\Phi(\mathbf{x}).
\end{equation}

Equation~\eqref{eq:newtonian-loop-potential} makes explicit
that the gravitational field depends on the spatial
distribution of the loop's energy, rather than solely on
its geometric asymmetry or its maximum radial extent.

\subsection{The Multipole Expansion}

Let
\[
r=|\mathbf{x}|,
\qquad
\mathbf{n}=\frac{\mathbf{x}}{r},
\]
and suppose the origin is chosen so that
\[
a_\gamma
\coloneq
\sup_{s\in S^1}|\mathbf{y}(s)|
\]
is finite. In the far-field regime $r\gg a_\gamma$, we may
expand the Newtonian kernel as
\begin{equation}
	\label{eq:kernel-multipole-expansion}
	\frac{1}{|\mathbf{x}-\mathbf{y}|}
	=
	\frac{1}{r}
	+\frac{\mathbf{n}\cdot\mathbf{y}}{r^2}
	+\frac{3(\mathbf{n}\cdot\mathbf{y})^2-|\mathbf{y}|^2}
	{2r^3}
	+O\!\left(\frac{|\mathbf{y}|^3}{r^4}\right).
\end{equation}
Substituting \eqref{eq:kernel-multipole-expansion} into
\eqref{eq:newtonian-loop-potential} yields the first three
multipole contributions:
\begin{equation}
	\label{eq:loop-multipole-potential}
	\Phi(\mathbf{x})
	=
	-G\left[
	\frac{M_{\mathrm{eff}}}{r}
	+\frac{\mathbf{n}\cdot\mathbf{d}}{r^2}
	+\frac{n_iQ_{ij}n_j}{2r^3}
	+O\!\left(
	\frac{M_{\mathrm{eff}}a_\gamma^3}{r^4}
	\right)
	\right].
\end{equation}
Here and below, repeated Cartesian indices are summed.

The moments appearing in \eqref{eq:loop-multipole-potential}
are defined as follows.

\paragraph{Monopole.}
The monopole is the total effective mass
\begin{equation}
	M_{\mathrm{eff}}
	=
	\oint_\gamma\lambda_{\mathrm{eff}}(s)\,ds.
\end{equation}
It determines the leading $r^{-1}$ contribution to the
potential and is insensitive to the detailed arrangement of
the energy along the loop.

\paragraph{Dipole.}
The mass-dipole vector is
\begin{equation}
	\label{eq:loop-dipole}
	\mathbf{d}
	=
	\oint_\gamma
	\lambda_{\mathrm{eff}}(s)\mathbf{y}(s)\,ds.
\end{equation}
It records the first spatial moment of the energy distribution
relative to the chosen origin. If the origin is placed at the
effective centre of mass,
\begin{equation}
	\mathbf{y}_{\mathrm{cm}}
	=
	\frac{1}{M_{\mathrm{eff}}}
	\oint_\gamma
	\lambda_{\mathrm{eff}}(s)\mathbf{y}(s)\,ds,
\end{equation}
then $\mathbf{d}=0$ by construction. Thus, a nonzero dipole
relative to the centre of the reference circle can reflect
a displacement of the energy centroid, rather than an
intrinsic shape asymmetry alone.

\paragraph{Quadrupole.}
The trace-free mass-quadrupole tensor is defined by
\begin{equation}
	\label{eq:loop-quadrupole}
	Q_{ij}
	=
	\oint_\gamma
	\lambda_{\mathrm{eff}}(s)
	\left(
	3y_i(s)y_j(s)
	-|\mathbf{y}(s)|^2\delta_{ij}
	\right)ds,
\end{equation}
where $\delta_{ij}$ is the Kronecker delta. By construction,
$Q_{ii}=0$. The quadrupole captures directional information
about the spatial distribution of the effective mass that
is absent from the monopole. Its contribution to the
potential decays as $r^{-3}$.

Higher multipoles describe progressively finer features of
the source distribution. Their contributions are suppressed
by additional powers of the ratio $a_\gamma/r$ in the
far-field expansion.

\subsection{Relation to Hemispheric Energy Asymmetry}

Recall the normalized energy asymmetry
$\mathcal{A}_E$ defined in
\eqref{eq:energy-asymmetry}. This quantity compares the
integrated energies of the two hemispheres, whereas the
multipole moments depend on the full spatial distribution
of energy. Consequently, $\mathcal{A}_E$ does not determine
the dipole or quadrupole by itself.

Nevertheless, the two descriptions provide complementary
information. The total energy determines the monopole in
our effective model, while changes in peak heights and
spacings alter the spatial weighting in
\eqref{eq:loop-dipole} and \eqref{eq:loop-quadrupole}.
We can therefore investigate whether variations of
$(A_1,N_1,A_2,N_2)$ produce systematic changes in the
multipole moments.

For this comparison, we retain the centre of the reference
unit circle as a fixed origin. We may additionally compute
the moments about the effective centre of mass to distinguish
a displacement of the energy centroid from changes in the
loop's intrinsic energy distribution.

\subsection*{Scope and Limitations}

The expansion \eqref{eq:loop-multipole-potential} is useful
when the observation point is far from the loop relative
to its overall size. If a brane lies close to the string,
the expansion may converge slowly or provide an inadequate
approximation; in that case, the potential
\eqref{eq:newtonian-loop-potential} should be evaluated
directly.

Moreover, this section concerns a quasi-static Newtonian
potential. Multipole expansions of gravitational radiation
involve time-dependent source moments and additional
relativistic considerations \cite{Thorne1980MultipoleRadiation}.
They are not interchangeable with the static potential
expansion used here.

Finally, the effective line-mass density
\eqref{eq:effective-line-mass} is a modeling assumption.
The present calculation does not establish that geometric
asymmetry creates gravity or that the gravitational force
is proportional to $\mathcal{A}_E$. Rather, it provides
a framework for testing how the energy distribution
predicted by our elastic-loop model affects the resulting
approximate gravitational field.

```latex
\section{Brane Deformations and Diffeomorphic Minimization}
\label{sec:brane-deformations}

Let \(X\) be a smooth manifold, let
\(\iota_0:\Sigma^n\hookrightarrow X\) be an initial brane embedding,
and let \(\mathcal F_{\mathfrak B}\) denote the brane's flow set (Def.~\ref{brane}).
Each admissible ambient diffeomorphism \(\Phi:X\to X\) determines
a deformed embedding
\[
\iota=\Phi\circ\iota_0.
\]
We write
\[
\mathcal O_{\mathfrak B}(\iota_0)
=
\left\{
\Phi\circ\iota_0
\;\middle|\;
\Phi \text{ is an admissible endpoint of a flow in }
\mathcal F_{\mathfrak B}
\right\}
\]
for the family of brane configurations reachable from
\(\iota_0\) through the prescribed deformations. Since these
deformations arise from ambient diffeomorphisms, they preserve the
brane's diffeomorphism type, keeping stable any potential Aharanov--Bohm-like properties, and remaining topologically stable.

Fix a smooth closed string
\(\gamma:S^1\hookrightarrow X\). The flow set determines which
brane configurations are admissible, but does not by itself specify
which configuration is preferred relative to \(\gamma\). To express
this dependence, let
\(\mathcal E_{\mathrm{br}}[\iota]\) denote the brane energy and
let \(\mathcal E_{\mathrm{int}}[\gamma,\iota]\) denote the
string--brane interaction energy. We seek a brane configuration
minimizing their sum over the admissible orbit:
\begin{equation} \label{Master}
	\boxed{
		\iota^\ast(\gamma)
		\in
		\underset{\iota\in\mathcal O_{\mathfrak B}(\iota_0)}
		{\operatorname{arg\,min}}
		\left(
		\mathcal E_{\mathrm{br}}[\iota]
		+
		\mathcal E_{\mathrm{int}}[\gamma,\iota]
		\right).
	}
\end{equation}
Here the string is held fixed while the brane configuration is
selected. The dependence of \(\iota^\ast\) on \(\gamma\) expresses
the possibility that different string geometries favor different
brane configurations. The interaction term is the mechanism through
which this geometric compatibility enters the minimization problem.

The formulation is conditional on the existence of a minimizer:
the admissible orbit need not be compact, and the energy need not
attain its infimum without additional hypotheses. If several
configurations attain the same minimum, the expression
\(\operatorname{arg\,min}\) denotes the full set of minimizers, and
\(\iota^\ast(\gamma)\) is any one of its elements. Thus the flow set
specifies the allowed deformations, whereas the energy functional
selects energetically preferred configurations whenever a minimum
exists.

This selection principle does not, by itself, prescribe a time
evolution. A dynamical model would additionally specify how the
string and brane respond to their combined energy, allowing their
geometries to evolve together rather than minimizing over the brane
with the string held fixed. Here is one possible example:
\begin{eg}[Gradient-flow coupling]
	Let
	\[
	\mathcal J[\gamma,\iota]
	=
	\mathcal E_{\mathrm{str}}[\gamma]
	+
	\mathcal E_{\mathrm{br}}[\iota]
	+
	\mathcal E_{\mathrm{int}}[\gamma,\iota]
	\]
	be the total energy of a string--brane configuration
	\(q=(\gamma,\iota)\). Choose a product metric on the configuration
	space, with metrics \(\mathsf G_{\mathrm{str}}\) and
	\(\mathsf G_{\mathrm{br}}\) on the string and brane components,
	respectively. One possible dynamical coupling is the negative
	gradient flow
	\[
	\begin{aligned}
		\partial_t\gamma_t
		&=-\operatorname{grad}_{\gamma}
		\mathcal J[\gamma_t,\iota_t],\\
		\partial_t\iota_t
		&=-\operatorname{grad}_{\mathcal O_{\mathfrak B},\iota}
		\mathcal J[\gamma_t,\iota_t],
	\end{aligned}
	\]
	where the brane gradient is taken along the admissible
	diffeomorphic orbit
	\(\mathcal O_{\mathfrak B}(\iota_0)\). Thus the brane evolves
	only through permitted deformations, while the string is free
	to change its shape.
	
	The interaction energy
	\(\mathcal E_{\mathrm{int}}[\gamma,\iota]\) couples the two
	evolution equations: changes in the string geometry affect the
	brane's energy landscape, and changes in the brane geometry
	affect the forces acting on the string. Whenever this gradient
	flow is well-defined, the total energy satisfies
	\[
	\frac{d}{dt}\mathcal J[\gamma_t,\iota_t]
	=
	-\|\operatorname{grad}_{\gamma}\mathcal J\|_{\mathsf G_{\mathrm{str}}}^{2}
	-\|\operatorname{grad}_{\mathcal O_{\mathfrak B},\iota}\mathcal J\|_{\mathsf G_{\mathrm{br}}}^{2}
	\leq 0.
	\]
	Hence the coupled evolution dissipates total energy, subject
	to the chosen metric and admissibility constraints. This is
	one possible variational model rather than a uniquely determined
	physical law.
\end{eg}

\begin{rk}
	The use of a gradient flow is motivated by established methods
	in geometric analysis and elasticity. Elastic curves have been
	studied through gradient flows of curvature-dependent energies
	\cite{DziukKuwertSchaetzle2002}, while diffeomorphic surface
	flows provide an energy-minimizing framework for topology-preserving
	surface evolution \cite{ZhangEtAl}. Coupled variational treatments
	of elastic curves and surfaces also demonstrate how interactions
	between geometric objects can be incorporated into a common energy
	\cite{Atai}. Motivated by these approaches, we propose the negative
	gradient flow of \(\mathcal J\) as one possible coupling of the
	string and brane. This choice defines a dissipative relaxation
	model; it is not asserted to be the unique or fundamental physical
	law governing their dynamics.
\end{rk}

\section{Self-Intersections and Controlled Extensions}
\label{sec:self-intersections}

The preceding analysis assumes that the string is a smooth embedded
closed curve \(\gamma:S^1\hookrightarrow X\). This assumption
provides a regular geometric setting for the elastic energy, the
induced gravitational field, and the coupled string--brane
variational problem. We now discuss how the framework may be
extended to admit controlled self-intersections.

The four-parameter planar model introduced above does not itself
admit self-intersections. Indeed, for a positive radial profile
\(r(\theta)>0\), the parametrization
\[
\gamma(\theta)
=
\bigl(r(\theta)\cos\theta,\,
r(\theta)\sin\theta,\,0\bigr)
\]
is injective on \(S^1\). Thus self-intersecting configurations
require an enlargement of the admissible class of string geometries.

For this purpose, we replace the embedding condition by the
assumption that \(\gamma:S^1\to X\) is a smooth immersion, so that
its derivative never vanishes, while allowing distinct parameter
values to map to the same spatial point. Define the set of
double-point pairs by
\[
\mathcal D(\gamma)
=
\left\{
\{s_1,s_2\}\subset S^1
\;\middle|\;
s_1\neq s_2,\quad
\gamma(s_1)=\gamma(s_2)
\right\}.
\]
As a first controlled extension, we restrict attention to immersions
for which \(\mathcal D(\gamma)\) is finite, each self-intersection
has exactly two preimages, and the two branches have distinct
tangent directions. In the planar setting, this means that the
branches cross with nonparallel tangents. Triple intersections,
tangential contacts, and singular points of the parametrized curve
are excluded at this stage.

\paragraph{Elastic energy.}
For such an immersion, the elastic energy remains well-defined
as an integral over the parameterized curve:
\[
\mathcal E_{\mathrm{str}}[\gamma]
=
\int_\gamma
\left(
\tau+\frac{\kappa_b}{2}\kappa^2
\right)\,ds.
\]
At a double point \(p=\gamma(s_1)=\gamma(s_2)\), curvature is
evaluated separately along the two parameterized branches. Since
an isolated crossing has zero arclength measure, it contributes
no additional point energy to this local functional. A separate
crossing interaction would have to be introduced explicitly if
the model is intended to assign energy to the intersection itself.

\paragraph{Gravitational field and multipoles.}
The effective line density and gravitational potential may likewise
be defined by integration over the parameterized curve:
\[
\Phi_\gamma(\mathbf{x})
=
-G\int_\gamma
\frac{\lambda_{\mathrm{eff}}(s)}
{|\mathbf{x}-\gamma(s)|}\,ds,
\qquad
\lambda_{\mathrm{eff}}(s)
=
\frac{\varepsilon_{\mathrm{str}}(s)}{c^2}.
\]
At a self-intersection, both branches contribute to the source
integral. The intersection does not, by itself, introduce an
additional point mass. The monopole and higher multipole moments
therefore remain definable by the same parameterwise integrals,
provided the effective density is specified along each branch.

A distinction is necessary when evaluating the field on the string
itself. The idealization of the string as a source of zero thickness
can produce singular self-field expressions, including near
self-intersections. A calculation of gravitational self-force
therefore requires an appropriate regularization or a formalism
designed to handle the string's own gravitational backreaction.
Such issues arise in the relativistic cosmic-string literature,
including work that treats loops with self-intersections
\cite{ChernoffFlanaganWardell2019StringBackreaction}.

\paragraph{Consequences for coupled evolution.}
Allowing self-intersections changes the admissible configuration
space for the string, but does not by itself change the definition
of the brane flow set \(\mathcal F_{\mathfrak B}\). The coupled
energy
\[
\mathcal J[\gamma,\iota]
=
\mathcal E_{\mathrm{str}}[\gamma]
+
\mathcal E_{\mathrm{br}}[\iota]
+
\mathcal E_{\mathrm{int}}[\gamma,\iota]
\]
can be considered on a suitable space of immersed loops and
admissible brane embeddings, provided its constituent terms are
well-defined there. A gradient-flow model would then describe
evolution within this enlarged configuration space, subject to
the regularity and existence assumptions required for the flow.

Geometric self-intersection must, however, be distinguished from
physical intercommutation. In cosmic-string models, intersecting
segments may reconnect, producing different loop configurations
or fragmenting a parent loop into daughter loops
\cite{CasperAllen1995RealisticLoops,Thompson1988CosmicStringDynamics,Vachaspati1989SelfIntersections}. Such a transition changes the
string's connectivity and cannot be specified by a smooth
single-loop gradient flow alone. To include it, the model would
need an additional reconnection rule and a prescription for
updating the configuration and its energy.

These distinctions are also important for interpreting gravitational
radiation. Self-intersecting loop trajectories have been studied
in calculations of gravitational radiation spectra
\cite{StormScherrer2024GVLoops}, while relativistic treatments
of gravitational backreaction examine the forces generated by
the string's own field
\cite{ChernoffFlanaganWardell2019StringBackreaction}. These works
provide relevant physical precedents, but they concern cosmic
strings and their gravitational radiation or backreaction; they
do not establish that the elastic-loop model proposed here is
itself a fundamental graviton model.

Accordingly, the results derived for embedded loops remain the
baseline case of this paper. Self-intersecting immersions provide
a controlled extension, while branch interactions, self-force
regularization, and topology-changing reconnection are additional
modeling questions to be addressed separately.

\appendix
```latex
\subsection{Numerical Example}
\label{subsec:numerical-example}

We now evaluate the elastic energy and gravitational multipole
moments for the asymmetric loop of
Figure~\ref{fig:asymmetric-closed-loop}. This example illustrates
how the shape parameters affect the energy distribution and the
resulting effective gravitational field.

\paragraph{Dimensionless energy.}
The reference circle in our geometric model has normalized radius
one. Let \(R_0\) denote its physical radius and introduce the
dimensionless bending parameter
\[
\beta=\frac{\kappa_b}{\tau R_0^2}.
\]
Writing \(\widehat{\kappa}=R_0\kappa\) and
\(d\widehat{s}=ds/R_0\), the dimensionless energy of hemisphere
\(I_i\) is
\[
\widehat{E}_i
\coloneq \frac{E_i}{\tau R_0}
=
\int_{I_i}
\left(1+\frac{\beta}{2}\widehat{\kappa}^2\right)
\sqrt{r^2+(r')^2}\,d\theta.
\]
For this illustrative calculation, we set \(\beta=1\).
This is a normalized modeling choice, not an empirically determined
material parameter.

We use the shape parameters
\[
(A_1,N_1,A_2,N_2)=(0.8,2,0.45,4),
\]
as in Figure~\ref{fig:asymmetric-closed-loop}. Numerical
integration gives the following dimensionless contributions.
Here \(\widehat{L}_i\) denotes the length contribution and
\[
\widehat{B}_i
=
\frac12\int_{I_i}\widehat{\kappa}^{\,2}\,d\widehat{s}
\]
denotes the bending contribution.

\begin{center}
	\begin{tabular}{lrrr}
		\toprule
		Hemisphere & \(\widehat{L}_i\) & \(\widehat{B}_i\)
		& \(\widehat{E}_i\) \\
		\midrule
		\(I_1\) & \(5.7180\) & \(17.2018\) & \(22.9198\) \\
		\(I_2\) & \(5.6733\) & \(84.0956\) & \(89.7689\) \\
		\midrule
		Total & \(11.3913\) & \(101.2974\) & \(112.6887\) \\
		\bottomrule
	\end{tabular}
\end{center}

Consequently, the normalized energy asymmetry is
\[
\mathcal{A}_E
=
\frac{\widehat{E}_2-\widehat{E}_1}
{\widehat{E}_1+\widehat{E}_2}
\approx 0.5932.
\]
Thus, approximately \(59.3\%\) is the signed normalized
difference between the two hemispheric energy contributions.
In this example, the second hemisphere contributes substantially
more energy despite having the smaller peak amplitude.

The length contributions are nearly equal; the difference is
dominated by the bending energy. This illustrates that peak
spacing, and not merely peak height, can strongly influence the
elastic energy distribution. The numerical result depends on the
chosen bending parameter \(\beta\) and should not be interpreted
as a universal relation between shape asymmetry and energy.

\paragraph{Effective mass and multipole moments.}
We next calculate the moments relative to the fixed origin at
the center of the reference circle. Let
\(\widehat{\mathbf y}=\mathbf y/R_0\) and define the
dimensionless energy-density factor
\[
\widehat{\varepsilon}(\theta)
=
1+\frac12\widehat{\kappa}(\theta)^2.
\]
The dimensionless monopole, dipole, and quadrupole are
\[
\widehat{M}
=
\oint_\gamma
\widehat{\varepsilon}\,d\widehat{s},
\qquad
\widehat{\mathbf d}
=
\oint_\gamma
\widehat{\varepsilon}\,
\widehat{\mathbf y}\,d\widehat{s},
\]
and
\[
\widehat{Q}_{ij}
=
\oint_\gamma
\widehat{\varepsilon}
\left(
3\widehat{y}_i\widehat{y}_j
-|\widehat{\mathbf y}|^2\delta_{ij}
\right)d\widehat{s}.
\]
The dimensional moments are recovered through
\[
M_{\mathrm{eff}}
=\frac{\tau R_0}{c^2}\widehat{M},
\qquad
\mathbf d
=\frac{\tau R_0^2}{c^2}\widehat{\mathbf d},
\qquad
Q_{ij}
=\frac{\tau R_0^3}{c^2}\widehat{Q}_{ij}.
\]
Because \(\beta=1\), the dimensionless monopole equals the
dimensionless total energy. Numerically, we obtain
\[
\widehat{M}\approx112.6887,
\qquad
\widehat{\mathbf d}
\approx
\begin{pmatrix}
	0\\
	-46.1093\\
	0
\end{pmatrix},
\]
and
\[
\widehat{Q}
\approx
\begin{pmatrix}
	79.7427 & 0 & 0\\
	0 & 79.7427 & 0\\
	0 & 0 & -159.4854
\end{pmatrix}.
\]
The quadrupole tensor is trace-free, as required by its definition.
The vanishing \(x\)-component of the dipole follows from the
reflection symmetry of this particular loop about the vertical
axis. Its negative \(y\)-component indicates that the effective
energy centroid lies below the reference origin. This is an
origin-dependent dipole, not an intrinsic asymmetry measure
independent of the chosen coordinates.

\paragraph{Far-field potential.}
As a final check, consider observation points
\(\mathbf{x}_{\pm}=(0,\pm10R_0,0)\). The maximum loop radius is
\(1.8R_0\), so these points lie outside the immediate source
region. Using the monopole, dipole, and quadrupole terms gives
\[
\begin{aligned}
	\Phi(\mathbf{x}_+)
	&\approx
	-\frac{G\tau}{c^2}
	\left(
	\frac{\widehat{M}}{10}
	+\frac{\widehat{d}_y}{10^2}
	+\frac{\widehat{Q}_{yy}}{2\cdot10^3}
	\right)\\
	&\approx
	-\frac{G\tau}{c^2}(10.8476),
	\\[6pt]
	\Phi(\mathbf{x}_-)
	&\approx
	-\frac{G\tau}{c^2}
	\left(
	\frac{\widehat{M}}{10}
	-\frac{\widehat{d}_y}{10^2}
	+\frac{\widehat{Q}_{yy}}{2\cdot10^3}
	\right)\\
	&\approx
	-\frac{G\tau}{c^2}(11.7698).
\end{aligned}
\]
Direct numerical evaluation of the potential integral
\eqref{eq:newtonian-loop-potential}, using the same normalized
source density, gives respectively
\[
\Phi(\mathbf{x}_+)
\approx-\frac{G\tau}{c^2}(10.8462),
\qquad
\Phi(\mathbf{x}_-)
\approx-\frac{G\tau}{c^2}(11.7715).
\]
The close agreement illustrates the accuracy of the first three
multipole terms for these observation points. The potential is
more negative on the negative \(y\)-axis, consistently with the
downward displacement of the effective energy centroid.

This example demonstrates, within the stated effective Newtonian
model, how a chosen loop geometry determines its elastic energy,
effective mass moments, and approximate far-field potential.
It does not establish that geometric asymmetry itself generates
gravity: the field calculation depends on the assumed effective
line-mass density, and the numerical results depend on the
dimensionless bending parameter.


\begin{thebibliography}{9}
\bibitem{Adams}
J.~F.~Adams. \emph{Infinite Loop Spaces}. Hermann Weyl Lectures, The Institute for Advanced Study (AM-90). (1978).

\bibitem{Atai}
A.~A.~Atai. and D.~J.~Steigmann \emph{Coupled Deformations of Elastic Curves and Surfaces}. \emph{International Journal of Solids and Structures}, Vol.~35, No.~16, pp.~1915--1952. (1998).

\bibitem{CasperAllen1995RealisticLoops}
P.~Casper and B.~Allen.
\emph{Gravitational Radiation from Realistic Cosmic String Loops}.
\emph{Physical Review D}, Vol.~52, No.~8, pp.~4337--4348.
\href{https://doi.org/10.1103/PhysRevD.52.4337}
{doi:10.1103/PhysRevD.52.4337}. (1995).

\bibitem{ChernoffFlanaganWardell2019StringBackreaction}
D.~F.~Chernoff, \'{E}.~\'{E}.~Flanagan, and B.~Wardell. \emph{Gravitational Backreaction on a Cosmic String: Formalism}. \emph{Physical Review D}, Vol.~99, 084036. \href{https://doi.org/10.1103/PhysRevD.99.084036}{doi:10.1103/PhysRevD.99.084036}. (2019).

\bibitem{DziukKuwertSchaetzle2002}
G.~Dziuk, E.~Kuwert, and R.~Sch\"atzle.
\emph{Evolution of Elastic Curves in $\mathbb{R}^n$: Existence and Computation}.
\emph{SIAM Journal on Mathematical Analysis}, Vol.~33, No.~5, pp.~1228--1245.
\href{https://doi.org/10.1137/S0036141001383709}{doi:10.1137/S0036141001383709}.
(2002).

\bibitem{Green}
M.~B.~Green, J.~H.~Schwarz, and E.~Witten. \emph{Superstring Theory (Vol.~1 \& 2)}. Cambridge University Press, Cambridge. (1987).

\bibitem{PoissonWill2014Gravity}
E.~Poisson and C.~M.~Will. \emph{Gravity: Newtonian, Post-Newtonian, Relativistic}. Cambridge University Press, Cambridge. \href{https://doi.org/10.1017/CBO9781139507486}{doi:10.1017/CBO9781139507486}. (2014).

\bibitem{Polchinski}
J.~Polchinski. \emph{String Theory (Vol.~1 \& 2)}. Cambridge University Press, Cambridge. (1998).

\bibitem{Schreiber}
U.~Schreiber. \emph{From Loop Space Mechanics to Nonabelian Strings}. PhD thesis. \href{https://arxiv.org/abs/hep-th/0509163}{arXiv:hep-th/0509163}. (2005).

\bibitem{StormScherrer2024GVLoops}
S.~D.~Storm and R.~J.~Scherrer.
\emph{Gravitational Radiation Power Spectrum of Garfinkle--Vachaspati Cosmic String Loops}.
\emph{Physical Review D}, Vol.~110, 083531.
\href{https://doi.org/10.1103/PhysRevD.110.083531}
{doi:10.1103/PhysRevD.110.083531}. (2024).

\bibitem{Thompson1988CosmicStringDynamics}
C.~Thompson. \emph{Dynamics of Cosmic String}.
\emph{Physical Review D}, Vol.~37, pp.~283--297.
\href{https://doi.org/10.1103/PhysRevD.37.283}
{doi:10.1103/PhysRevD.37.283}. (1988).

\bibitem{Thorne1980MultipoleRadiation}
K.~S.~Thorne. \emph{Multipole Expansions of Gravitational Radiation}. \emph{Reviews of Modern Physics}, Vol.~52, No.~2, pp.~299--339. \href{https://doi.org/10.1103/RevModPhys.52.299}{doi:10.1103/RevModPhys.52.299}. (1980).

\bibitem{Vachaspati1989SelfIntersections}
T.~Vachaspati. \emph{Cosmic-String Loop Self-Intersections and Intercommuting}.
\emph{Physical Review D}, Vol.~39, pp.~1768--1771.
\href{https://doi.org/10.1103/PhysRevD.39.1768}
{doi:10.1103/PhysRevD.39.1768}. (1989).

\bibitem{ZhangEtAl}
S.~Zhang, L.~Younes, J.~Zweck, and J.~T.~Ratnanather. \emph{Diffeomorphic Surface Flows: A Novel Method of Surface Evolution}. \emph{SIAM Journal on Applied Mathematics}, Vol.~68, No.~3, pp.~806--824. (2008).

\bibitem{Zwiebach}
B.~Zwiebach. \emph{A First Course in String Theory}. Cambridge University Press, Cambridge. (2004).

\end{thebibliography}


\end{document}
